\section{결론}
\label{sec:conclusion}

저자들은 멀티 소스 뱅킹 사용자 이력을 위한 인코더형 파운데이션 모델 계열인 PRAGMA를 제시했다.
PRAGMA는 키--값--시간 토큰화 방식을, 프로필 상태와 이벤트를 각각 처리하는 두 인코더 가지(branch) 그리고 그 출력을 융합하는 히스토리 인코더와 결합하며, 대규모 이질적 금융 레코드 위에서 마스킹 모델링으로 사전학습된다.
다양한 다운스트림 과제 — 신용 평가, 사기 탐지, 커뮤니케이션 인게이지먼트, 상품 추천, 정기 거래 탐지, 라이프타임 밸류 예측 등 — 에 걸쳐, 단일 사전학습 백본이 원시 뱅킹 이벤트 시퀀스로부터 곧바로 우수한 성능을 달성하며, 금융 응용을 위한 범용 표현 계층을 제공한다.

저자들의 실험은 몇 가지 실용적 통찰을 드러낸다.
LoRA 파인튜닝은 일관되게 처음부터의 전체 학습과 맞먹거나 그 이상의 성능을 내면서도 갱신되는 파라미터는 일부에 그치며, 이는 사전학습 표현이 과제 간에 효과적으로 전이됨을 확인해 준다.
10\,M에서 1\,B 파라미터로의 스케일링은 신용 평가 같은 더 어려운 과제에서 큰 이득을 가져오는 반면, 라이프타임 밸류 예측 같은 과제에서는 더 작은 모델만으로도 경쟁력 있는 표현을 제공하므로 실용적인 효율성 트레이드오프가 가능하다.
전용 프로필 상태 인코더는 신용 평가나 사기 탐지처럼 정적 맥락 속성이 변별적인 과제에 특히 유용하며, 그런 신호가 덜 관련된 과제에서는 아키텍처가 우아하게 성능을 유지한다.
또한 사전학습된 텍스트 인코더의 통합은 텍스트 비중이 큰 도메인의 성능을 끌어올리지만, 텍스트가 희박한 과제에서는 학습 오버헤드를 정당화하지 못하는 것으로 나타났다.
마지막으로, AML 사례 연구는 한 가지 분명한 한계를 부각한다 — 레코드 간 관계 구조에 의존하는 과제는, 이벤트 이력을 고립된 단위로 처리하는 모델로는 도달하기 어렵다.

이 결과들은, 이질적 구조(heterogeneous structure), 불규칙한 시간성(irregular timing), 운영적 제약(operational constraint)에도 불구하고 멀티 소스 뱅킹 이벤트 시퀀스가 텍스트나 비전과 같은 방식으로 \emph{전이 가능한 표현}을 허용함을 시사한다.
자금세탁방지 같은 관계적 과제를 위해 모델을 레코드 간 상호작용까지 포착하도록 확장하는 것은 향후 연구의 유망한 방향이다.


\subsubsection*{감사의 글}

저자들은 이 연구에 기여해 준
Dmitry Mittov,
Ian Iakobsen,
Aleksandr Pushin,
Muhammad Anas,
Viacheslav Karpov,
Nathalie Skrzypek,
Leyla Sultanova,
Francisco Sanz Estevez,
Nikita Kravchuk,
Tadas Krisciunas,
Amey Baokar,
Hanna Danilovich,
Jyoti Prakash Bal,
Vitalii Radchenko,
Kade Main,
Nic Hatia,
그리고 그 밖의 Revoluter 동료들에게 감사를 전한다.
